% =============================================================================
% Kapitel 1: Einleitung
% =============================================================================
\chapter{Einleitung}
\label{ch:einleitung}

Die Elektromobilität ist ein immer wichtiger werdendes Thema. Sie trägt zur Reduzierung von Treibhausgasemissionen bei, verringert die Abhängigkeit zu fossilen Brennstoffen und somit auch wirtschaftliche Abhängigkeiten. Gleichzeitig bietet sie eine neue Möglichkeit bei der Energieversorgung.

Die Möglichkeit bei der Energieversorgung ist hier besonders interessant. Fahrzeuge, die mit Strom betrieben werden, können nicht nur als Verbraucher betrachtet werden, sondern auch als Energiespeicher. Elektrofahrzeuge können also nicht nur Energie aufnehmen, sondern diese auch wieder abgeben, was besonders im Kontext von erneuerbaren Energien von großer Bedeutung ist.

Im Moment bestehen noch kaum rentable Möglichkeiten, überschüssige Energie aus erneuerbaren Quellen zu speichern, weshalb dem Thema eine besondere Relevanz zukommt: Durch eine zuverlässige Speicherung könnte Energie auch dann genutzt werden, wenn Solar- oder Windanlagen nicht aktiv erzeugen. Damit ließe sich der Anteil erneuerbarer Energien an der Versorgung deutlich erhöhen, was ein wichtiger Beitrag zur Energiewende wäre.

Jedoch bringt eine Speicherung von Energie in Elektrofahrzeugen einen gerinfügigen Mehrwert - ökonomisch wie auch ökologisch - wenn diese nicht zuverlässig gesteuert wird. Es wird ein System benötigt, das die intelligente Steuerung von Energieflüssen zwischen Fahrzeug und Stromnetz ermöglicht. Eine intelligente Steuerung ist nur möglich, wenn das Nutzerverhalten identifiziert werden kann. Es muss also ein System entworfen werden, das eine zuverlässige Vorhersage über das Mobilitätsverhalten des Nutzers ermöglicht.

% -----------------------------------------------------------------------------
\section{Motivation}
\label{sec:motivation}

In dieser Arbeit wird untersucht, inwiefern sich das Nutzungsverhalten von Elektrofahrzeugen vorhersagen lässt, um  potentiell den Energiefluss zwischen Fahrzeug und Stromnetz zu optimieren. Voraussetzung hierfür ist die Verfügbarkeit von V2G (Vehicle-to-Grid), also der Fähigkeit von Elektrofahrzeugen, Energie nicht nur aufzunehmen, sondern auch wieder ans Stromnetz abzugeben. Obwohl V2G zunehmend an Präsenz gewinnt, bestehen in der praktischen Anwendung weiterhin große Lücken. Ein zentraler Baustein für das effiziente Zusammenwirken von Elektrofahrzeugen und Stromnetz ist ein zuverlässiges Vorhersagemodell des Mobilitätsverhaltens: Es erlaubt, Lade- und Entladezyklen so zu steuern, dass sowohl die Effizienz der mobilen Akkumulatoren als auch die Auslastung des Netzes optimiert werden. Entscheidend ist dabei die Zuverlässigkeit der Vorhersage, um ein fehlerloses System gewährleisten zu können.
% -----------------------------------------------------------------------------
\section{Problemstellung}
\label{sec:problemstellung}

Die größte Hürde stellt eine stabile Datengrundlage dar. Da der Anwendungsfall sehr anspruchsvoll und dynamisch ist, lässt er sich mit bestehenden Ansätzen bislang nur unzureichend abbilden. Typischerweise erzielen Machine Learning Modelle sehr gute Ergebnisse bei dieser Art von Vorhersageproblemen. Voraussetzung dafür ist allerdings eine hochqualitative Datenbasis, denn Fehleinschätzungen können zu einer fehlerhaften Energieverwaltung und im schlimmsten Fall zur Nichtverfügbarkeit des Fahrzeugs führen. Es müssen also verlässliche Daten über das Mobilitätsverhalten erhoben oder von Dritten bezogen werden. Da Qualität in diesem Bereich nicht genau definiert ist, gilt es dies an bestimmten Kriterien und Werten festzulegen.

In Betrachtung der existierenden wissenschaftlichen Literatur besteht im Moment ein Mangel an Studien, die sich explizit mit dem Mobilitätsverhalten von Individuen und der Erstellung von Prognosemodellen für diesen Zweck beschäftigen. Es gibt wertvolle Studien, die sich mit dem Mobilitätsverhalten von Gruppen beschäftigen, jedoch liefern diese wenig Aufschluss darüber, wie sich der Individualverkehr vorhersagen lässt. Damit künftig eine intelligente Steuerung von Energieflüssen möglich wird, müssen zuverlässige Vorhersagemodelle entwickelt werden, die auf hochqualitativen Daten zum individuellen Mobilitätsverhalten basieren. 

Bisherige Arbeiten zur Prognose im Mobilitäts- und Energiekontext konzentrieren sich überwiegend auf den Ladezeitpunkt, die Leistung oder den letztendlichen Ladezustand sowie auf aggregiertes Nutzungsverhalten ganzer Gruppen. Das individuelle Fahr- und Nutzungsverhalten einzelner Fahrzeuge wird dabei kaum modelliert, obwohl gerade dieses für eine genaue Steuerung bzw. Vorhersage von Nöten ist. Wo individuelles Fahrverhalten betrachtet wird, ist es zwischen den Fahrern stark heterogen, sodass ein einzelnes, übergreifendes Modell nicht den kompletten Kontext erfässt; zugleich ist die Datenbasis der wenigen einschlägigen Studien häufig nicht ausreichend oder -- wie bei synthetisch erzeugten Profilen -- nicht zuverlässig repräsentativ für reale Mobilitätsdaten. Eine ausführliche Einordnung der relevanten Forschungsstränge und die daraus abgeleitete Forschungslücke finden sich in \Cref{sec:stand_forschung}.

% -----------------------------------------------------------------------------
\section{Zielsetzung}
\label{sec:zielsetzung}

Das Ziel der Arbeit ist es, zuverlässige Modelle zur Prognose für individuelles Mobilitätsverhalten zu entwickeln. Die Modelle sollen den Zweck erfüllen, genauer Uhrzeiten vorhersagen zu können, wann ein Fahrzeug fährt und wann es parkt. Beispielsweise soll ein Modell, für bis zu sieben Tage,  Dazu sind die erhobenen Daten, qualitativ zu bewerten und wenn nötig, aus Datenschutzgründen zu anonymisieren. Auf dieser Grundlage werden Modelle erstellt, die das Mobilitätsverhalten von Nutzern möglichst genau vorhersagen können. Des Weiteren soll Gegenstand der Arbeit sein, die Vorhersagegenauigkeit verschiedener Modelle zu vergleichen und deren Einflussfaktoren zu quantifizieren. Die Ergebnisse sollen dazu beitragen, eine intelligente Steuerung von Energieflüssen zwischen Fahrzeug und Stromnetz zu ermöglichen, um so die Nutzung von erneuerbaren Energien zu verbessern und eine wichtige Rolle bei der Energiewende zu spielen.

% -----------------------------------------------------------------------------
\section{Forschungsfragen}
\label{sec:forschungsfragen}

\begin{enumerate}
    \item Die Beschaffung und Analyse von verlässlichen Daten über das Mobilitätsverhalten von Nutzern, um so die Vorhersage zu ermöglichen.
    \begin{enumerate}
        \item Berücksichtigung von Datenschutzaspekten bei der Datenerhebung.
    \end{enumerate}
    \item Erstellung mehrerer Modells zur Vorhersage des Mobilitätsverhaltens von Nutzern.
    \begin{enumerate}
        \item Vergleich der Modelle hinsichtlich ihrer Vorhersagegenauigkeit.
        \item Identifikation der wichtigsten Einflussfaktoren auf das Mobilitätsverhalten.
    \end{enumerate}
\end{enumerate} 

% -----------------------------------------------------------------------------
\section{Aufbau der Arbeit}
\label{sec:aufbau}

\todo{Kurze Beschreibung der Struktur der Arbeit. Welches Kapitel behandelt welches Thema?}

Die vorliegende Arbeit ist wie folgt aufgebaut:

\textbf{\Cref{ch:grundlagen}} erläutert die theoretischen Grundlagen, die für das Verständnis dieser Arbeit notwendig sind.

\textbf{\Cref{ch:methodik}} beschreibt die gewählte Methodik und das Vorgehen.

\textbf{\Cref{ch:implementierung}} stellt die Implementierung der entwickelten Lösung vor.

\textbf{\Cref{ch:evaluation}} präsentiert die Evaluation und die erzielten Ergebnisse.

\textbf{\Cref{ch:diskussion}} diskutiert die Ergebnisse und deren Bedeutung.

\textbf{\Cref{ch:fazit}} fasst die Arbeit zusammen und gibt einen Ausblick auf zukünftige Arbeiten.
